% Base acadêmica: decisões e fontes em ../ABNT.md.
% Compilar com LuaLaTeX. Documento em redação, não versão de entrega.
\documentclass[12pt,oneside,openany,a4paper,english,brazil]{abntex2}
\usepackage{fontspec}
\IfFontExistsTF{Times New Roman}{
  \setmainfont{Times New Roman}
}{
  \setmainfont{TeX Gyre Termes}
  \PackageWarningNoLine{vaulty}{Times New Roman indisponivel; usando TeX Gyre Termes}
}
\usepackage{graphicx}
\usepackage[a4paper,left=3cm,right=2cm,top=3cm,bottom=2cm,
  headheight=12pt,headsep=\dimexpr1cm-12pt\relax]{geometry}

% Escolhas do modelo: 12 pontos no texto; 10 nas exceções.
\renewcommand{\ABNTEXchapterfont}{\normalfont\bfseries}
\renewcommand{\ABNTEXchapterfontsize}{\normalsize}
\renewcommand{\ABNTEXsectionfont}{\normalfont\bfseries}
\renewcommand{\ABNTEXsectionfontsize}{\normalsize}
\renewcommand{\ABNTEXsubsectionfont}{\normalfont\itshape}
\renewcommand{\ABNTEXsubsectionfontsize}{\normalsize}
\renewcommand{\ABNTEXfontereduzida}{\fontsize{10}{12}\selectfont}
\renewcommand{\footnotesize}{\fontsize{10}{12}\selectfont}
\setlength{\ABNTEXcitacaorecuo}{4cm}
\setlength{\parindent}{1.25cm}
\setlength{\parskip}{0pt}
\setsecnumdepth{subsubsection}
\settocdepth{subsubsection}
\setlength{\beforechapskip}{0pt}
\setlength{\afterchapskip}{\baselineskip}
\setbeforesecskip{\baselineskip}
\setaftersecskip{\baselineskip}
\setbeforesubsecskip{\baselineskip}
\setaftersubsecskip{\baselineskip}
\captionnamefont{\normalfont\normalsize}
\captiontitlefont{\normalfont\normalsize}
\captiondelim{ --- }
\captionstyle{\SingleSpacing\raggedright}
\setlength{\footmarkwidth}{0.6em}
\setlength{\footmarksep}{0pt}
\renewcommand{\footnoterule}{\kern-3pt\hrule width 5cm\kern2.6pt}
\hypersetup{hidelinks,pdftitle={Vaulty: documentação do projeto integrador}}

% TODO: substituir pelos dados oficiais, sem atribuir autoria não confirmada.
\autor{[Nome completo do integrante 1]\\
  [Nome completo do integrante 2]\\
  [Nome completo do integrante 3]\\
  [Nome completo do integrante 4]}
\titulo{VAULTY: CENTRALIZAÇÃO DE INFORMAÇÕES OPERACIONAIS}
\instituicao{Centro Universitário SATC\\[Nome oficial do curso]}
\local{[Cidade de apresentação]}
\data{2026}
\orientador{[Nome e titulação do orientador]}
\preambulo{Trabalho de Projeto Integrador apresentado ao
  [nome oficial do curso] do Centro Universitário SATC, como requisito
  para [objetivo acadêmico definido pela disciplina].}

\newcommand{\pendente}[1]{\textit{[A preencher: #1.]}}
\newcommand{\natureza}{%
  \noindent\hspace*{0.5\textwidth}%
  \begin{minipage}{0.5\textwidth}
    \SingleSpacing\imprimirpreambulo
  \end{minipage}\par}
\newcommand{\fonte}[1]{\par\begingroup
  \ABNTEXfontereduzida\SingleSpacing\noindent Fonte: #1\par\endgroup}
\newcommand{\referencia}[1]{\noindent #1\par\vspace{\baselineskip}}

% Estilo próprio: número no alto à direita também na abertura das seções.
\makepagestyle{vaulty}
\makeoddhead{vaulty}{}{}{\ABNTEXfontereduzida\thepage}
\makeevenhead{vaulty}{}{}{\ABNTEXfontereduzida\thepage}
\makeoddfoot{vaulty}{}{}{}
\makeevenfoot{vaulty}{}{}{}

\newif\iffichacatalografica
\fichacatalograficafalse % Habilitar somente com os dados da biblioteca.

\begin{document}
\selectlanguage{brazil}
\frenchspacing
\OnehalfSpacing
\pretextual

% Capa fora da contagem; composição em 12 pontos.
\begin{capa}
  \centering
  {\imprimirinstituicao\par}
  \vspace{2cm}
  {\imprimirautor\par}
  \vfill
  {\bfseries\imprimirtitulo\par}
  \vfill
  {\imprimirlocal\par}
  {\imprimirdata\par}
\end{capa}

% Folha de rosto: primeira página contada, sem número visível.
\setcounter{page}{1}
\begin{titlingpage}
  \thispagestyle{empty}
  \begin{center}
    \imprimirautor\par
    \vspace{4cm}
    {\bfseries\imprimirtitulo\par}
  \end{center}
  \vspace{2cm}
  \natureza
  \vspace{1cm}
  \noindent\hspace*{0.5\textwidth}%
  \begin{minipage}{0.5\textwidth}
    \SingleSpacing Orientador: \imprimirorientador
  \end{minipage}
  \vfill
  \begin{center}\imprimirlocal\par\imprimirdata\end{center}
\end{titlingpage}
\setcounter{page}{2}

\iffichacatalografica
  \thispagestyle{empty}
  \vspace*{\fill}
  % TODO: inserir ficha fornecida pela biblioteca, por exemplo como imagem.
  \pendente{dados internacionais de catalogação fornecidos pela biblioteca}
  \clearpage
  \addtocounter{page}{-1} % Página catalográfica não integra a contagem.
\fi

% Folha de aprovação sem título; não declara aprovação já ocorrida.
\begin{folhadeaprovacao}
  \begin{center}
    \imprimirautor\par
    \vspace{2cm}
    {\bfseries\imprimirtitulo\par}
  \end{center}
  \vspace{1cm}
  \natureza
  \vspace{1cm}
  \noindent Data de aprovação: \pendente{preencher após aprovação}
  \vspace{1cm}
  \assinatura{[Nome, titulação e instituição do examinador 1]}
  \assinatura{[Nome, titulação e instituição do examinador 2]}
  \assinatura{[Nome, titulação e instituição do examinador 3]}
\end{folhadeaprovacao}

\begin{resumo}
  \noindent\pendente{resumo final em parágrafo único, com objetivo,
  método, resultados comprovados e conclusão; extensão recomendada
  de 150 a 500 palavras}
  \vspace{\baselineskip}
  \noindent\textbf{Palavras-chave:} engenharia de requisitos;
  informação operacional; controle de acesso.
\end{resumo}

\begin{resumo}[Abstract]
  \begin{otherlanguage*}{english}
    \noindent\pendente{tradução fiel do resumo final para o inglês}
    \vspace{\baselineskip}
    \noindent\textbf{Keywords:} requirements engineering;
    operational information; access control.
  \end{otherlanguage*}
\end{resumo}

% Ativar listas somente quando houver figuras/tabelas reais.
% \listoffigures*\cleardoublepage
% \listoftables*\cleardoublepage
\tableofcontents*\cleardoublepage

\textual
\pagestyle{vaulty}
\aliaspagestyle{chapter}{vaulty}

\chapter{INTRODUÇÃO}
O Vaulty é um sistema proposto para centralizar informações operacionais,
como chaves de API e documentação, em um ambiente com controle de acesso.
O problema definido pelo proponente é o tempo gasto procurando essas
informações em conversas pessoais, além da dificuldade de disponibilizá-las
a outros desenvolvedores que também necessitam delas.

Este trabalho apresenta a concepção e a organização dos requisitos do
projeto integrador. A documentação distingue as decisões confirmadas pelo
proponente, a revisão coletiva ainda pendente e as evidências de implementação
e verificação. O benefício de reduzir o tempo de procura é esperado e ainda
precisa de avaliação; não se apresenta como resultado medido.

\section{Objetivo geral}
Especificar uma solução que permita aos usuários autorizados encontrar e
consultar informações operacionais cadastradas, sem depender de uma nova
solicitação pessoal para localizar cada informação.

\section{Objetivos específicos}
Definir os tipos de conteúdo e seu ciclo de vida; modelar as ações permitidas
e o acesso aos recursos; estabelecer critérios de aceite e condições de
verificação; relacionar requisitos às tarefas de desenvolvimento; e registrar
decisões e mudanças de escopo.

\section{Organização do documento}
A apresentação adota a estrutura de trabalho acadêmico da Associação
Brasileira de Normas Técnicas (ABNT), conforme a NBR 14724:2024
(ABNT, 2024). As seções seguintes tratam da fundamentação, do método,
dos requisitos, da verificação e das considerações finais.

\chapter{FUNDAMENTAÇÃO TEÓRICA}
\pendente{conceitos de Engenharia de Requisitos utilizados no projeto,
com referências bibliográficas efetivamente consultadas}

\chapter{MÉTODO E CONCEPÇÃO}
\section{Origem e elaboração dos requisitos}
O projeto teve origem na proposta do proponente. A especificação foi refinada
por discussão das necessidades, regras de acesso, fluxos de uso e critérios de
aceite. O registro documental de 12 de outubro de 2026 identifica a versão
candidata 0.2, confirmada pelo proponente e pendente de revisão dos demais
integrantes. A validação externa dos requisitos permanece pendente.

\section{Público e escopo}
O público contempla desenvolvedores e profissionais de infraestrutura.
A primeira entrega abrange textos, links, segredos e documentos com um
arquivo por item. A criação de chaves de API permanece no processo externo
ao produto. Vídeos e plataformas Apple não integram a entrega definida.

\section{Equipe e restrição de prazo}
A equipe possui quatro integrantes, dois no backend e dois no frontend,
com disponibilidade concentrada nos finais de semana. O prazo informado
para apresentação à banca é 27 de novembro de 2026. A estimativa do esforço
restante e a revisão do compromisso de entrega pelo grupo estão pendentes.

\chapter{ESPECIFICAÇÃO DOS REQUISITOS}
\section{Organização e rastreabilidade}
A especificação atual identifica 53 requisitos funcionais (RF), 17 requisitos
não funcionais (RNF) e critérios de aceite (CA). Os casos de uso (UC)
descrevem os fluxos consolidados. Os identificadores devem ser preservados
na transposição para este documento e na ligação com tarefas e testes.

\section{Síntese do acesso ao conteúdo}
O papel de um membro determina as ações permitidas; grupos e
compartilhamentos determinam quais conteúdos estão acessíveis. A operação
exige autorização nos dois aspectos. Pastas permitem acesso habitual por
grupos, enquanto itens podem conceder acesso individual como exceção.
Busca, listagem e consulta devem respeitar o acesso efetivo ao conteúdo.

\section{Complementação da especificação}
\pendente{transpor os requisitos, regras de negócio e critérios de aceite
de REQUISITOS.md; os cinco casos de uso de CASOS-DE-USO.md;
o modelo de MODELO-CONCEITUAL.md; e a matriz de RASTREABILIDADE.md}

\chapter{PLANO DE VERIFICAÇÃO E RESULTADOS}
Os critérios documentados definem o comportamento a verificar. O plano
de desempenho contempla seis usuários em contextos isolados e 120 buscas;
a execução e suas evidências ainda precisam ser registradas. A recuperação
do backup também depende de verificação em ambiente descartável.

\pendente{descrever ambiente real, procedimentos, resultados e limitações
de cada cenário efetivamente executado, preservando falhas e requisitos
ainda não atendidos}

\chapter{CONSIDERAÇÕES FINAIS}
\pendente{avaliar os objetivos à luz das evidências, registrar limitações
e trabalho futuro; não antecipar conclusões de testes pendentes}

\postextual
\chapter*{REFERÊNCIAS}
\addcontentsline{toc}{chapter}{REFERÊNCIAS}
\begin{SingleSpace}
  \referencia{ASSOCIAÇÃO BRASILEIRA DE NORMAS TÉCNICAS.
    \textbf{ABNT NBR 14724:2024}: informação e documentação:
    trabalhos acadêmicos: apresentação. Rio de Janeiro: ABNT, 2024.
    Disponível em:
    \url{https://funesa.se.gov.br/wp-content/uploads/2025/06/TRABALHOS-ACADEMICOS.pdf}.
    Acesso em: 12 out. 2026.}
\end{SingleSpace}

% Apêndices: material produzido pelos autores. Anexos: material de terceiros.
% Ativar quando existir conteúdo; manter numeração contínua e fontes dos anexos.
% \begin{apendicesenv}
%   \chapter{INSTRUMENTOS DE VERIFICAÇÃO}
% \end{apendicesenv}
\end{document}
